% Generated draft: requires mathematical and source review.
\documentclass[11pt]{article}
\usepackage[margin=25mm]{geometry}
\usepackage{fontspec}
\setmainfont{Latin Modern Roman}
\usepackage{amsmath,amssymb,mathtools,unicode-math}
\setmathfont{Latin Modern Math}
\usepackage{enumitem,needspace,xurl,hyperref}
\setlength{\parindent}{0pt}
\setlength{\parskip}{5pt}
\setlength{\emergencystretch}{3em}
\newcommand{\N}{\mathbb{N}}
\newcommand{\Z}{\mathbb{Z}}
\newcommand{\Q}{\mathbb{Q}}
\newcommand{\R}{\mathbb{R}}
\newcommand{\C}{\mathbb{C}}
\newcommand{\F}{\mathbb{F}}
\newcommand{\A}{\mathbb{A}}
\newcommand{\PP}{\mathbb{P}}
\newcommand{\E}{\mathbb{E}}
\newcommand{\eps}{\varepsilon}
\newcommand{\dd}{\,\mathrm{d}}
\begin{document}
\section*{UnsolvedMath ID: 800835}
\textbf{Source problem code:} catalogue entry
\subsection{Definitions and mathematical statement}
The supplied excerpt concerns a system of balls and refers to families of balls denoted by $F_N$, with a limiting procedure $N\to\infty$. It does not define the underlying system, the families $F_N$, their equations of motion, their initial or boundary conditions, or the energy under consideration. Consequently, these objects cannot be assigned precise mathematical definitions from the supplied record alone.

The excerpt takes as its starting point that energy has escaped from the system. It describes a proposed ``limit absorption principle'': introduce friction acting on the balls of $F_N$ so that this friction absorbs energy, and then let $N\to\infty$. The source objects that this procedure installs an energy sink from the outset and merely moves it farther and farther away; it therefore regards this explanation of energy escape as circular.

\textbf{Problem 34 (source question).} For the ball system referenced in the excerpt, is there a more natural way to introduce energy dissipation without explicitly introducing a friction mechanism?

\textbf{Unresolved formulation.} The supplied record does not specify what counts as a more natural construction, which energy should dissipate, how energy escape is measured, or which limiting or observational procedure is intended. In particular, it does not distinguish loss of energy from an observed subsystem from loss of total energy. No precise existence statement or dissipation inequality, with specified hypotheses and quantifiers, can be extracted without the missing system definitions and a criterion for the requested dissipation.
\subsection{Short English statement}
The question asks how energy dissipation could arise in the referenced system of balls without explicitly adding friction. The source regards friction placed progressively farther away as a circular explanation of energy escape. The underlying dynamics and the intended meaning of dissipation are missing, so this remains a qualified research question requiring mathematical review.
\subsection{Sources}
\begin{itemize}
\item Local UnsolvedMath registry, permanent ID \texttt{800835}, source problem code \texttt{catalogue entry}, title: \emph{Geometric Hydrodynamics in Open Problems --- Problem 34}. The supplied registry attributes the problem to Khesin, Misiolek, and Modin, with proposed year 2022, and labels it open. This label is supplied source metadata, not an independently verified assessment.
\item Khesin, Misiolek, and Modin, \emph{Geometric Hydrodynamics in Open Problems} (2022), Problem 34. Supplied source URL: \url{https://arxiv.org/abs/2205.01143}. The supplied excerpt is the basis of this statement; the cited document has not been independently consulted here.
\item Supplied provenance records access on \texttt{2026-10-16} and extraction by \texttt{arxiv\_latex}. These are supplied provenance claims, not a report of independent access or verification. The registry records creation and update at \texttt{2026-10-16T14:00:00Z} and upstream revision \texttt{d371ee917ac9647fbb06cd6fc529deec6f3dda82}.
\end{itemize}
\end{document}
